\begin{table}[H]

\centering
\caption{Lista de programas a cargo del club ambiental}
\label{tab:porgramas_del_club_ambiental}

\begin{threeparttable}

{\small\setstretch{1.25}\rowcolors{2}{Grey}{White}

\begin{tabular}{M{0.15\linewidth}M{0.375\linewidth}M{0.375\linewidth}}\hline
    
    \theader{Programa} & \theader{Funcionamiento} & \theader{Objetivos} \\ \hline

    Compostaje & Compostaje de los residuos orgánicos universitarios y de comercios cercanos, realizado por alumnos voluntarios y utilizado como abono en áreas verdes universitarias y prácticas académicas.
    &
    \begin{itemize}
        \item Crear una herramienta más para el aprendizaje de los alumnos
        \item Fomentar el manejo de residuos responsable
    \end{itemize} \\

    Difusión & Actividades dentro y fuera de la universidad para dar difusión a temas de cuidado ambiental.
    &
    \begin{itemize}
        \item Crear conciencia de la importancia del cuidado ambiental
        \item Dar herramientas a los estudiantes para llevar una vida más ecológica
        \item Dar a conocer avances en tecnologías de cuidado ambiental
    \end{itemize} \\

    Limpiezas & Eventos que convoquen a la comunidad estudiantil para realizar limpieza, tanto dentro como fuera de la universidad.
    &
    \begin{itemize}
        \item Fomentar la disposición responsable de residuos
        \item Mantener una universidad libre de contaminación
        \item Recolectar materia compostable
    \end{itemize} \\

    Vinculación & Contacto con organizaciones locales para asistir a sus eventos como representantes universitarios y dar difusión de estos entre la comunidad estudiantil.
    &
    \begin{itemize}
        \item Dar apoyo a organizaciones que fomenten el cuidado ambiental
        \item Reforestación (a futuro)
        \item Cultivo de árboles fertilizados con el programa de compostaje y su uso para proyectos de reforestación universitarios o municipales
        \item Dar apoyo a la formación de un municipio más sustentable
    \end{itemize} \\

    \hline
\end{tabular}

}

\begin{tablenotes}\tiny % ex: \tnote{1} ... \item[1]
    \item[1] 
\end{tablenotes}

\end{threeparttable}
\end{table}